
\section{Overview: ML as a Paradigm Shift}
\label{sec:ml-overview}

Machine learning is the idea that complex rules or behaviors can be learned directly from data by optimization algorithms, rather than programmed explicitly by human experts.
In classical data analysis, a domain expert encodes their knowledge as a fixed set of rules that are then applied to data to extract a quantitative summary of the information.
It is worth emphasizing that this approach often works very well.
Particle physicists have become very adept at choosing relevant features of the data and learning from them.
However, this approach is rarely optimal, especially in cases where the optimal solution exists but is too complex to specify by hand.
To draw an example from \Cref{ch:tagging}, no expert-designed feature set captures the correct decision-surface for classifying jets initiated by boosted top quarks from the light quark and gluon jet background.
This is due to a lack of knowledge about certain features of the jet formation process (e.g. hadronization), but also just due to complexity.
Specifying the optimal feature by hand is probably not possible, but machine learning methods allow this complexity to be learned from the data.

Machine learning methods organize naturally into several paradigms based on the structure of the available training information:

\begin{itemize}
    \item \textit{Supervised learning} trains a model using labeled examples, where each input is paired with a target output.
    Classification and regression are the canonical supervised tasks, and the most widely used in particle physics.
    \Cref{ch:tagging} is entirely dedicated to the jet tagging classification task.
    \item \textit{Unsupervised learning} operates without labels, with the goal of learning structure in the input data. Generative models that learn to produce samples from $p_{\text{data}}$ without access to class labels are a prominent example and are the focus of much of \cref{ch:unfolding}.
    \item \textit{Self-supervised learning} is a hybrid approach in which the supervisory signal is derived from the data itself rather than from external annotation. This is most commonly used to learn general-purpose representations rather than to solve a specific downstream task. In HEP, self-supervised pretraining underlies some recent foundation models for jet physics, which use techniques such as masked particle modeling to learn representations of jet constituents that transfer across tasks~\cite{Golling:2024abg,Birk:2024knn}.
    \item \textit{Reinforcement learning} (RL) trains an agent to maximize a reward signal through interaction with an environment. The key feature of RL is that it does not require a differentiable training objective, though some methods construct a surrogate in practice. RL is responsible for landmark results in game playing~\cite{Silver:2016alphago} and is the basis for the post-training pipelines of modern large language models~\cite{Ouyang:2022instructgpt}. Direct applications of RL within HEP have so far been limited~\cite{Baretz:2025zsv}, though its relevance is growing alongside the agentic AI systems discussed in \cref{ch:conclusion}.
\end{itemize}

Machine learning is not a new idea in particle physics.
Boosted decision trees (BDTs) have been in use in HEP since 2004, and were applied to the single top quark discovery at the Tevatron~\cite{D0:2006ngk,CDF:2009itk}.\footnote{The history of BDTs in HEP is interesting.
The original application of BDTs illustrated that they \textit{outperformed} artificial neural networks.
Following the deep learning revolution initiated by AlexNet, BDTs were largely supplanted by neural networks, which proved far more expressive on the high-dimensional data common in HEP.
More recently, BDTs have seen something of a resurgence, driven by their remarkable training speed relative to neural networks, and their robustness against noisy or irrelevant features in the density ratio estimation tasks covered in \cref{sec:dre}~\cite{Finke:2023ltw}.}
What changed with the deep learning era was not the concept of learned classifiers but the scale of the models.
Neural networks scale better than BDTs, and so could handle the more complicated tasks that have driven much of the interest in ML in recent years.

The success of machine learning in HEP rests on two structural advantages.
The first is access to high-quality simulated data produced by the Monte Carlo event generators discussed in \cref{sec:jets}.
These simulated datasets carry \textit{truth level} information that can be used to derive effective training labels.
This is a significant advantage over other application domains where manual labeling is a significant bottleneck.
Label assignment in HEP has its own subtleties, which are discussed in the context of top tagging in \Cref{ch:tagging}, but effective labels can often be set using straightforward generator-level rules.
The second advantage is that HEP data is generated by physical processes that are well understood and are known to respect specific symmetries.
Lorentz invariance, permutation invariance of jet constituents, and translational equivariance in the $(\eta, \phi)$ plane are all properties that have been used to aid neural network optimization.
These properties could equally well be learned from the data and the community is starting to explore this direction, but exploiting them was a major benefit to early deep learning applications in HEP.

Given that deep neural networks are the dominant paradigm in machine learning at present, the remainder of this chapter provides an overview of the core network architectures and training methods used throughout this thesis.
